\begin{document}
\docTitle{Results guidelines}
All experimental results included in a document must be generated programmatically.
This includes tables, figures, statistical tests, and any other result artifacts.

\section*{Artifact generation}
Each result artifact must be generated by a Python script that takes exactly one JSON file as input:

\begin{DocCode}
python3 script_name.py data.json
\end{DocCode}

Here, \code{script_name.py} and \code{data.json} are placeholders.
Each artifact must identify the actual script and JSON input paths used to generate it.
The same script or JSON input file may be used for multiple artifacts.

\section*{Data collection}
Each input JSON file must be generated programmatically by a Python script named:

\begin{DocCode}
python3 collect_<json_name>.py
\end{DocCode}

The collection script must gather data from the relevant source locations and produce the JSON consumed by the artifact-generation scripts.
It must not compute statistics or aggregate results.

The resulting JSON must be a list of entries.
Each entry must contain:
\begin{itemize}
    \item fields that uniquely identify the experiment;
    \item the location of the original data source;
    \item the commit hash associated with the data;
    \item the script used to obtain or process the data.
\end{itemize}

\section*{Provenance}
Every generated result artifact must include provenance information.
\agentProposedEdit{agent-remove-43a76bc3-b943-5885-aff5-2420298dd0a0}{inconsistent with authoritative tex closure}{For each generated experimental table, create a corresponding provenance \code{<table_file_name>_provenance.json}.}{}
\agentProposedEdit
  {agent-add-99eb6fac-d789-549d-a720-207ae984bc2e}
  {Separate compact reproducibility metadata from bulky generator payloads while preserving artifact identity, as prepared in Codex conversation hash 01a03d38-a300-7be1-ac03-4b540a4f117d.}
  {}
  {Compact tracked provenance belongs under \path{scripts/provenance/}, while raw result rows, remote snapshots, and bulky generation payloads belong under the ignored \path{.cache/generation/} tree.
The tracked \path{scripts/provenance/artifact_manifest.json} records each compiled artifact's relative path, SHA-256 digest, source location, generator identity, and corresponding cached or external payload identity.
The \path{tex/generated/} tree contains only derived files consumed by the current canonical PDF and never contains sibling provenance JSON or raw result collections.}

The provenance JSON must contain one entry for every result value or table cell.
Each entry must include:
\begin{itemize}
    \item the row identifier in the generated artifact;
    \item the column identifier in the generated artifact;
    \item the LaTeX value shown in the artifact;
    \item the unique identifier of the source row used to produce the value;
    \item the original source-row data required to reproduce or verify the value.
\end{itemize}


\section*{Examples}
\paragraph{Tables}
Generated experimental tables must be stored under:

\begin{DocCode}
tex/generated/tables/<document_or_topic>/<table_name>.tex
\end{DocCode}

They must be imported directly from that path.
For example:

\begin{DocCode}
\begin{table}[ht]
\centering
\caption{Caption text}
\label{tab:acc_main_experiment_table}
\input{tex/generated/tables/main/acc_main_experiment_table.tex}
\end{table}
\end{DocCode}

Every generated experimental table file must end with:

\begin{DocCode}
% This table was generated using python3 script_name.py data.json
% data.json was created using python3 collect_<json_name>.py
% Provenance of the data can be found at <table_file_name>_provenance.json
\end{DocCode}

These comments must contain the actual artifact-generation script, input JSON file, collection script, and provenance file used for the table.

\section*{Generated figures}

Generated figures must likewise be included directly from their generated file path.
For example:

\begin{DocCode}
\begin{figure}[H]
  \centering
  \includegraphics[width=0.95\linewidth]{tex/generated/plots/results/accuracy_vs_cost_balldrop.pdf}
  \caption{Episode top-1 accuracy versus mean evaluation cost for \texttt{BallDrop} under the four main-experiment conditions. Markers and error bars report the mean and sample standard deviation over five seeds, and the dashed horizontal line marks random-chance accuracy.}
  \label{fig:accuracy-vs-cost-balldrop}
\end{figure}
\end{DocCode}

\section*{Important requirements}

\begin{itemize}
    \item \agentImportant{Do not use \texttt{\textbackslash IfFileExists}, \texttt{\textbackslash InputIfFileExists}, or similar conditional file-existence logic. Generated artifacts must be referenced directly.}
    \item \agentImportant{Every table and figure must have a caption. 
    Captions must be written directly in the document source and must not be generated programmatically. 
    They are short (around three sentences or so at most). 
    Do not contains results.
    They provide an overview or description of how the data were generated}
\end{itemize}

\end{document}
